% Part: first-order-logic
% Chapter: sequent-calculus
% Section: provability-quantifiers

\documentclass[../../../include/open-logic-section]{subfiles}

\begin{document}

\olfileid{fol}{seq}{qpr}
\olsection{\usetoken{S}{derivability} आणि संख्यापक}

\begin{explain}
  संपूर्णता प्रमेयासाठी या विभागात स्थापित केलेली~$\Proves$ संबंधी तथ्ये
  क्रमवर्ती कलनातील नियमांनी निष्पन्न होतात, हेही आवश्यक आहे.
\end{explain}

\begin{thm}
\ollabel{thm:strong-generalization} जर $c$ हे $\Gamma$ किंवा $!A(x)$ मध्ये
न येणारे स्थिर असेल आणि $\Gamma \Proves !A(c)$, तर $\Gamma
\Proves \lforall[x][!A(x)]$.
\end{thm}

\begin{proof}
$\pi_0$ ही $\Gamma_0 \Sequent !A(c)$ या क्रमवर्तीची $\Log{LK}$-!!{derivation}
असलेली एखादी सांत $\Gamma_0 \subseteq \Gamma$ घ्या.  $\RightR{\lforall}$
अनुमान जोडल्यावर आपल्याला $\Gamma_0 \Sequent
\lforall[x][!A(x)]$ ची !!a{derivation} मिळते, कारण $c$ हे $\Gamma$ किंवा
$!A(x)$ मध्ये येत नाही आणि म्हणून आयगेन चराची अट पूर्ण होते.
\end{proof}

\begin{prop}
\ollabel{prop:provability-quantifiers}
\begin{tagenumerate}{prvEx,prvAll}
\tagitem{prvEx}{$!A(t) \Proves \lexists[x][!A(x)]$.}{}
\tagitem{prvAll}{$\lforall[x][!A(x)] \Proves !A(t)$.}{}
\end{tagenumerate}
\end{prop}

\begin{proof}
\begin{tagenumerate}{prvEx,prvAll}
\tagitem{prvEx}{$!A(t) \Sequent \lexists[x][!A(x)]$ ही क्रमवर्ती
  !!{derivable} आहे:
  \begin{prooftree}
    \Axiom$!A(t) \fCenter !A(t)$
    \RightLabel{\RightR{\lexists}}
    \UnaryInf$!A(t) \fCenter \lexists[x][!A(x)]$
    \end{prooftree}}{}
\tagitem{prvAll}{$\lforall[x][!A(x)] \Sequent !A(t)$ ही क्रमवर्ती
  !!{derivable} आहे:
  \begin{prooftree}
    \Axiom$!A(t) \fCenter !A(t)$
    \RightLabel{\LeftR{\lforall}}
    \UnaryInf$\lforall[x][!A(x)] \fCenter !A(t)$
    \end{prooftree}}{}
\end{tagenumerate}
\end{proof}

\end{document}
